% !TeX program = lualatex
\documentclass[11pt]{report} % Clase de documento para un informe
% Documento y codificación general
\usepackage[utf8]{inputenc} % Tildes, "ñ" y símbolos especiales.
% Idioma y traducción
\usepackage[spanish,es-tabla]{babel} % Traduce el documento al español.
\usepackage[a4paper,left=3cm,right=2.5cm,top=2.5cm,bottom=2.5cm]{geometry} % Permite cambiar el tamaño de la página.
% Imágenes y figuras
\usepackage{graphicx} % Permite incluir imágenes.
\usepackage{subcaption} % Permite agregar subfiguras dentro de una figura principal, cada una con su leyenda.
\usepackage{float} % Permite fijar la posición de las figuras y tablas [H].
% Matemáticas y símbolos
\usepackage{amsmath} % Permite usar entornos matemáticos avanzados.
\usepackage{amssymb} % Permite usar símbolos matemáticos adicionales.
\usepackage{mathtools} % Permite usar herramientas matemáticas avanzadas.
% Hipervinculos y referencias
\usepackage[hidelinks]{hyperref} % Permite crear hipervínculos transparentes dentro del documento.
\usepackage[table,xcdraw]{xcolor} % Permite usar colores textos y tablas.
\usepackage{setspace} % Permite ajustar el interlineado.
\usepackage{siunitx}

\begin{document}
    \begin{spacing}{1.5}
        \begin{flushleft}

                \begin{center}
                    \Large
                    \textbf{FORMULARIO DE AVANCE MENSUAL DE\\
                    PROYECTO INTEGRADOR}
                \end{center}
                \vspace{2em}

                FECHA: 27/05/2026

                ALUMNO: Tardón, Damián David

                TEMA: Diseño, desarrollo e implementación de un software para la adquisición y análisis de ondas de impulso tipo rayo (\SI[parse-numbers=false]{1,2/50}{\micro\second}).

                DIRECTOR: Blasco, Marcos Javier

                CO-DIRECTOR: Serra, Gabriel Horacio.

                FECHA ESTIMADA DE FINALIZACIÓN: Julio 2026.

                \hrulefill

                INFORME:

        \end{flushleft}
    \end{spacing}

    Se está redactando el informe del trabajo. Se tiene una primera versión del marco teórico, y se está empezando a redactar el marco metodológico.

\end{document}